\documentclass[11pt,letterpaper]{article}
\usepackage[T1]{fontenc}
\usepackage{lmodern,amsmath,amssymb,amsthm,microtype}
\usepackage[margin=1in]{geometry}
\usepackage[hidelinks]{hyperref}
\newtheorem{theorem}{Theorem}
\newcommand{\lastcase}{29}
\newcommand{\newbound}{30}
\title{Signature moments in the one-isolate height-six forcing case}
\author{Aafi Khan Malek\\ \small Independent researcher, Ahmedabad, India\\
\small\href{mailto:aafimalek2032@gmail.com}{aafimalek2032@gmail.com}}
\date{28 November 2026}
\begin{document}
\maketitle

\begin{abstract}
For a ground-minimal weak height-six counterexample with exactly one
isolated point and at least one forcing edge, the earlier pair-block
analysis verifies at least 14 forcing stars. Applying the height-three
Sauer bound after quotienting by occurrence signatures yields a smaller
cap on every two-leaf avoidance fiber. Weighted moments of the link
graphs, combined with exact finite arithmetic, then exclude every star
count from 14 through \lastcase. The new fiber inequality is checked
in Lean, while the final numerical exclusions are not yet formalized
end to end. The general height-six problem is still open in this work.
\end{abstract}

\begin{theorem}
In the one-isolate forcing case specified above there are at least
\newbound\ forcing stars.
\end{theorem}

\section{Sharper bounds for a pair fiber}

The starting structural facts are from the previous
\href{https://github.com/Aafimalek/novel-discoveries/blob/research/parallel-foundations/frankl-height-six/star-counting.md}{pair-block analysis}.
There are $k$ stars with exactly one leaf each and one isolated point.
Every pair of centers belongs to exactly one omitted-star block.
Distinct blocks intersect in at most one point. For $p,q$ in block $B$,
the fiber avoiding their two leaves has at least $3k+W(p,q)$ members,
where $W(p,q)$ counts blocks containing neither center.

Write $b=|B|$ and $h=k-b$. Join two outsiders if their unique block
meets $B$; let $E(B)$ be the number of these links and $c_B$ the number
of connected components. Closure makes the endpoints of each link
have identical occurrence signatures in the compressed pair fiber.
Retaining one coordinate per component and the isolate embeds that
fiber into an inclusion order on at most $c_B+1$ coordinates, and
preserves unions. The compressed family has no chain of four, so the
height-three Sauer bound gives
\[
3k+W(p,q)\le H(c_B+1),\qquad H(n)=1+n+\binom n2 .
\]
The Lean proof uses the number of distinct signatures, which is no
larger than $c_B$.

Suppose all blocks have size at most $r$. Each of the $k-2$ outsiders
of $p,q$ has at least $k-3-2(r-2)$ partners whose block avoids both.
Since one such block covers at most $r-1$ partners, putting
$d=\max(0,\lceil(k-3-2(r-2))/(r-1)\rceil)$ gives
$W\ge\lceil(k-2)d/r\rceil$. From any lower bound $w$ for some pair
in $B$, obtain $c_B\ge c$ from $H(c+1)\ge3k+w$. If $c>h$ the
block is impossible. Otherwise $E(B)\le e=\binom{h-c+1}{2}$ by
maximizing the clique edge capacity of the components.

\newpage
For each $a\in B$ the groups of outsiders in other blocks through
$a$ number at least $h-E_a$, where $E_a$ is the number of links
they contribute. Choose $p,q$ so that the remaining $b-2$ link
counts sum to at most $\lfloor(b-2)E(B)/b\rfloor$. Blocks
disjoint from $B$ account for the outside pairs not among its
links. These facts give a new lower bound
\[
W(p,q)\ge (b-2)h-\lfloor (b-2)e/b\rfloor
  +\left\lceil\frac{\binom h2-e}{\binom r2}\right\rceil .
\]
Iterating gives a cap $L_b$ on the link count, or excludes size $b$.

\section{Block counts and weighted moments}

Let $n_s$ count blocks of size $s\ge3$, let
$\kappa_s=\binom{s}{2}-1$, and put $\sigma=\sum_s n_s\kappa_s$.
There are at most $\binom{k}{2}$ pairs in these blocks.
For any $p,q\in B$, counting pairs outside $p,q$ gives
\[
W(p,q)\ge \binom{k-2}{2}-\binom{b-2}{2}+\kappa_b-\sigma .
\]
Indeed a block counted in $W$ contributes one plus its saving
$\binom{s}{2}-1$ pairs, while a block touching exactly one of $p,q$
contributes at most its saving. The unique block $B$ containing both
contributes $\binom{b-2}{2}$. This additional lower bound feeds back
into the iteration and makes every cap depend on the actual candidate
saving $\sigma$.

Put $\alpha_s=\binom{s-1}{2}$. For a selected collection $J$ of
block sizes, let $d_v$ be the sum of $\alpha_s$ over its blocks
through $v$. Distinct intersecting blocks meet once; a block of size
$s$ contributes $\alpha_s$ links to any block it intersects.
Expanding the sum of squares therefore gives
\[
D=\sum_v d_v=\sum_{s\in J}s\alpha_s n_s,\qquad
\sum_v d_v^2\le U=\sum_{s\in J}n_s(s\alpha_s^2+\alpha_s L_s).
\]
The minimum sum of $k$ integer squares at total $D$ occurs when the
scores differ by at most one. A count vector violating that bound
cannot occur.

For the whole collection set $d_s=s\alpha_s$ and
$u_s=s\alpha_s^2+\alpha_s L_s$. With at least one size-$r$ block,
the inequality $0\ge D^2-kU\ge 2zD-z^2-kU$ gives a compact
test for each $(r,\sigma)$ row. For an integer $z$, if
\[
(2zd_r-ku_r-z^2)\kappa_s
  +(\sigma-\kappa_r)(2zd_s-ku_s)>0
\]
for every admissible $s$, the whole row is impossible: extract one
size-$r$ block and bound the contribution per unit of remaining
saving. The implementation checks this inequality directly over
the integers. Rows without such a witness go to count enumeration.

\newpage
\section{The finite replay}

A local refinement is useful near saturation. The number of size-two
blocks is $M=\binom{k}{2}-\sum_s n_s\binom{s}{2}$. A vertex in a
chosen size-$b$ block has $k-b$ partners outside that block. Another
block of size $s$ through it uses $s-1$ partners and contributes
$\alpha_s$ links. At most $M$ size-two incidences can occur among
the chosen vertices, since no size-two block lies inside the chosen
block. Integer knapsack lower bounds on the link count, at each
local deficit and then across the chosen vertices, can therefore
reject further count vectors. They need not realize compatible
degree profiles: ignoring that condition only lowers the bound.

The replay enumerates possible largest block sizes and the entire
range of savings allowed by pair capacity. It verifies the integer
tangent inequalities when available, otherwise enumerates every
residual count vector allowed by pair capacity and savings. It
checks the weighted moments and pair deficits, and may additionally
check capacities on the support left after removing one size class.
No candidate survives for $14\le k\le\lastcase$. The all-pair design
is separately ruled out by
$3k+\binom{k-2}{2}>H(k-1)$. Together with the previous exclusion
through 13, this proves the stated restriction.

\paragraph{Reproducibility and scope.}
The accompanying plain Python scripts replay all finite inequalities
without an external solver. In Lean 4.36.0, the signature projection
and its application to the actual pair fiber compile with the
standard axioms only. The block count arithmetic and its assembly
into the new bound are not yet an end-to-end Lean theorem. The
previous fully Lean-checked numerical bound remains 14.

\paragraph{Contributions.}
Aafi Khan Malek selected and directed the research. ChatGPT assisted
with proof development, checking, formalization and writing. Codex
was not used. This note has not been peer reviewed.
\end{document}
